\documentclass[10pt,a4paper]{amsart}
\usepackage{amsmath,amssymb,amsthm,graphicx}
% Retained minimal preamble for the Stacks Project chapter more-morphisms.tex (More on Morphisms).
% Verbatim excerpt of the chapter preamble: Xy-pic setup (native preamble lines 14-18) and the native
% theorem declarations with the mathematical macros (native preamble lines 154-217). Copied mechanically;
% nothing added, removed or redefined. The wrapper class and the amsmath/amssymb/amsthm/graphicx base
% packages are supplied by the assembler, so the class/externaldocument part of the native preamble is omitted.
% For commutative diagrams we use Xy-pic
\usepackage[all]{xy}

% We use 2cell for 2-commutative diagrams.
\xyoption{2cell}
% Theorem environments.
%
\theoremstyle{plain}
\newtheorem{theorem}[subsection]{Theorem}
\newtheorem{proposition}[subsection]{Proposition}
\newtheorem{lemma}[subsection]{Lemma}

\theoremstyle{definition}
\newtheorem{definition}[subsection]{Definition}
\newtheorem{example}[subsection]{Example}
\newtheorem{exercise}[subsection]{Exercise}
\newtheorem{situation}[subsection]{Situation}

\theoremstyle{remark}
\newtheorem{remark}[subsection]{Remark}
\newtheorem{remarks}[subsection]{Remarks}

\numberwithin{equation}{subsection}

% Macros
%
\def\lim{\mathop{\mathrm{lim}}\nolimits}
\def\colim{\mathop{\mathrm{colim}}\nolimits}
\def\Spec{\mathop{\mathrm{Spec}}}
\def\Hom{\mathop{\mathrm{Hom}}\nolimits}
\def\Ext{\mathop{\mathrm{Ext}}\nolimits}
\def\SheafHom{\mathop{\mathcal{H}\!\mathit{om}}\nolimits}
\def\SheafExt{\mathop{\mathcal{E}\!\mathit{xt}}\nolimits}
\def\Sch{\mathit{Sch}}
\def\Mor{\mathop{\mathrm{Mor}}\nolimits}
\def\Ob{\mathop{\mathrm{Ob}}\nolimits}
\def\Sh{\mathop{\mathit{Sh}}\nolimits}
\def\NL{\mathop{N\!L}\nolimits}
\def\CH{\mathop{\mathrm{CH}}\nolimits}
\def\proetale{{pro\text{-}\acute{e}tale}}
\def\etale{{\acute{e}tale}}
\def\QCoh{\mathit{QCoh}}
\def\Ker{\mathop{\mathrm{Ker}}}
\def\Im{\mathop{\mathrm{Im}}}
\def\Coker{\mathop{\mathrm{Coker}}}
\def\Coim{\mathop{\mathrm{Coim}}}

% Boxtimes
%
\DeclareMathSymbol{\boxtimes}{\mathbin}{AMSa}{"02}

%
% Macros for moduli stacks/spaces
%
\def\QCohstack{\mathcal{QC}\!\mathit{oh}}
\def\Cohstack{\mathcal{C}\!\mathit{oh}}
\def\Spacesstack{\mathcal{S}\!\mathit{paces}}
\def\Quotfunctor{\mathrm{Quot}}
\def\Hilbfunctor{\mathrm{Hilb}}
\def\Curvesstack{\mathcal{C}\!\mathit{urves}}
\def\Polarizedstack{\mathcal{P}\!\mathit{olarized}}
\def\Complexesstack{\mathcal{C}\!\mathit{omplexes}}
% \Pic is the operator that assigns to X its picard group, usage \Pic(X)
% \Picardstack_{X/B} denotes the Picard stack of X over B
% \Picardfunctor_{X/B} denotes the Picard functor of X over B
\def\Pic{\mathop{\mathrm{Pic}}\nolimits}
\def\Picardstack{\mathcal{P}\!\mathit{ic}}
\def\Picardfunctor{\mathrm{Pic}}
\def\Deformationcategory{\mathcal{D}\!\mathit{ef}}
\usepackage{cleveref}
\crefname{originalreference}{original reference}{original references}
\Crefname{originalreference}{Original reference}{Original references}
\makeatletter
\newlabel{cohomology-lemma-finite-cohomology}{{08DN}{source}{}{sourcepointer-1}{}}
\newlabel{cohomology-lemma-finite-cohomology@cref}{{[originalreference][0][]08DN}{[1][0][]source}{}{}{}}
\makeatother
\crefname{originalreference}{original reference}{original references}
\Crefname{originalreference}{Original reference}{Original references}
\makeatletter
\newlabel{modules-lemma-finite-presentation-quasi-coherent}{{01BO}{source}{}{sourcepointer-2}{}}
\newlabel{modules-lemma-finite-presentation-quasi-coherent@cref}{{[originalreference][0][]01BO}{[1][0][]source}{}{}{}}
\makeatother
\crefname{originalreference}{original reference}{original references}
\Crefname{originalreference}{Original reference}{Original references}
\makeatletter
\newlabel{morphisms-lemma-i-star-equivalence}{{01QY}{source}{}{sourcepointer-3}{}}
\newlabel{morphisms-lemma-i-star-equivalence@cref}{{[originalreference][0][]01QY}{[1][0][]source}{}{}{}}
\makeatother
\title{Relative pseudo-coherence in low degrees for modules}
\author{The Stacks Project Authors}
\date{stacks-project:more-morphisms.tex:local-snapshot-20261011; GNU FDL 1.2 or later}
\begin{document}
\maketitle
\noindent\textit{Editorial note.} Counted span is the author's own proof environment, source lines 17052-17067, immediately following the statement at 17039-17050.\par\medskip
\noindent\textit{Editorial note.} Three cross-chapter pointers (Cohomology, Modules, Morphisms) are recorded in assembly.external\_references.\par\medskip
\noindent\textit{Editorial note.} The author's closing sentence 'The proof of (2) is similar but less involved.' is retained verbatim.\par\medskip
\noindent\textit{Allowed standard prerequisites (editorial).}\begin{itemize}
\item Pseudo-coherence of a module over a polynomial algebra. Cohomology, lemma cohomology-lemma-finite-cohomology, cited at proof line 17060. Identifies (-1)-pseudo-coherence of the pushed-forward module with finite presentation over A\textasciicircum{}n\_S.
\item Finitely presented modules are quasi-coherent. Modules, lemma modules-lemma-finite-presentation-quasi-coherent, cited at proof lines 17061-17062. Supplies one direction of the quasi-coherence clause.
\item Quasi-coherence is preserved and reflected by a closed immersion. Morphisms, lemma morphisms-lemma-i-star-equivalence, cited at proof lines 17063-17064. Transfers quasi-coherence between F and i\_*F, which closes part (3).
\item Definition of m-pseudo-coherence relative to the base. used at proof line 17053 for part (1). Part (1) is immediate from the definition, as the author states.
\end{itemize}\medskip
\ifdefined\hypertarget\hypertarget{sourcepointer-1}{}\fi
\noindent\textit{Source reference.} Original reference 08DN: focused-seeds/geometry/S9-morph-rel-n-pseudo-module/ref-tags.txt, source lines 5--5. Navigation pointer only: the counted proof cites this statement by \textbackslash{}ref as a standing input; the deduction the counted proof has to make is carried out inside the counted span, and no part of that deduction is delegated to this pointer.\par\medskip
\ifdefined\hypertarget\hypertarget{sourcepointer-2}{}\fi
\noindent\textit{Source reference.} Original reference 01BO: focused-seeds/geometry/S9-morph-rel-n-pseudo-module/ref-tags.txt, source lines 6--6. Navigation pointer only: the counted proof cites this statement by \textbackslash{}ref as a standing input; the deduction the counted proof has to make is carried out inside the counted span, and no part of that deduction is delegated to this pointer.\par\medskip
\ifdefined\hypertarget\hypertarget{sourcepointer-3}{}\fi
\noindent\textit{Source reference.} Original reference 01QY: focused-seeds/geometry/S9-morph-rel-n-pseudo-module/ref-tags.txt, source lines 7--7. Navigation pointer only: the counted proof cites this statement by \textbackslash{}ref as a standing input; the deduction the counted proof has to make is carried out inside the counted span, and no part of that deduction is delegated to this pointer.\par\medskip
\begin{lemma}
\label{lemma-rel-n-pseudo-module}
Let $X \to S$ be a morphism of schemes which is locally of finite type.
Let $\mathcal{F}$ be an $\mathcal{O}_X$-module. Then
\begin{enumerate}
\item $\mathcal{F}$ is $m$-pseudo-coherent relative to $S$ for all $m > 0$,
\item $\mathcal{F}$ is $0$-pseudo-coherent relative to $S$ if and only if
$\mathcal{F}$ is a finite type $\mathcal{O}_X$-module,
\item $\mathcal{F}$ is $(-1)$-pseudo-coherent relative to $S$ if and only if
$\mathcal{F}$ is quasi-coherent and finitely presented relative to $S$.
\end{enumerate}
\end{lemma}

\begin{proof}
Part (1) is immediate from the definition. To see part (3)
we may work locally on $X$ (both properties are local). Thus we
may assume $X$ and $S$ are affine. Choose a closed immersion
$i : X \to \mathbf{A}^n_S$. Then we see that $\mathcal{F}$ is
$(-1)$-pseudo-coherent relative to $S$ if and only if $i_*\mathcal{F}$
is $(-1)$-pseudo-coherent, which is true if and only if $i_*\mathcal{F}$
is an $\mathcal{O}_{\mathbf{A}^n_S}$-module of finite presentation, see
Cohomology, Lemma \ref{cohomology-lemma-finite-cohomology}.
A module of finite presentation is quasi-coherent, see
Modules, Lemma \ref{modules-lemma-finite-presentation-quasi-coherent}.
By Morphisms, Lemma \ref{morphisms-lemma-i-star-equivalence}
we see that $\mathcal{F}$ is quasi-coherent if and only if $i_*\mathcal{F}$
is quasi-coherent. Having said this part (3) follows. The proof of (2)
is similar but less involved.
\end{proof}

\end{document}
